\documentclass[a4paper,10pt]{article}
\usepackage[margin=18mm]{geometry}
\usepackage{graphicx}
\title{Report 1: EPC2361 SPICE Characterization}
\author{LossCalculation-SPB}
\begin{document}
\maketitle
\section{Purpose and scope}
This report characterizes the supplied EPC2361 transistor model, separately
from the SPB inverter efficiency and cooling report. It provides DC
temperature-dependent on-resistance and current-dependent switching energies,
voltage/current slopes and voltage stress. These are simulation results for
one device in each position, not measured device guarantees or a parallel-module model.
\section{DC on-resistance}
LTspice sweeps drain current with a 5 V gate bias at each fixed temperature.
The extracted quantity is $R_{DS(on)}=V_{DS}/I_D$, including model drain/source
access resistance. Zero current is excluded from the division. No self-heating
is solved. The 150 C point is a characterization endpoint at the absolute
temperature limit, not the module's design target.
\begin{center}\begin{tabular}{rr}
Temperature (C) & Model on-resistance at 50 A (mOhm) \\
\hline
@RDS@
\end{tabular}\end{center}
The model's typical fit can differ from the datasheet's maximum resistance and
its latest typical curves. It must not be silently substituted for a worst-case
conduction-loss calculation.
\par\noindent\includegraphics[width=\linewidth]{../../results/spice_characterization/rds_on.pdf}
\clearpage
\section{Switching characterization setup}
The previously completed double-pulse sweep supplies @CASES@ cases, of which
@FAILURES@ fail at least one modeled voltage screen. Its original configuration
is saved with this characterization; a model hash check prevents combining DC
results with a DPT table generated from a different extracted model.
\begin{center}\begin{tabular}{ll}
@SETUP@
\end{tabular}\end{center}
Internal gate resistance is additionally 0.4 ohm. The complementary gate stays
off. A first pulse builds inductor current; first turn-off and second turn-on
are measured. Horizontal axes use actual inductor current at the respective
edge, not the nominal current setting or the instantaneous ringing peak.
The DC temperature sweep does not imply a temperature sweep of switching energies.

\subsection{Energy and slope definitions}
Turn-on plots show excess DUT channel/access-resistor heat: integration of
observed internal dissipative power minus ordinary on-state conduction over
the measurement window. The conduction reference is a separate DC run of the
same model at the same temperature and 5 V gate bias. Turn-off retains delayed
channel conduction. These quantities exclude gate-drive loss and complementary
device reverse conduction. They are not standardized manufacturer $E_{on}/E_{off}$
specifications. Signed drain-terminal integrals are also retained in the CSV;
they differ from heat because stored capacitance energy changes during an edge.
Do not add Coss loss again to energies that already include its dissipation.

Both turn-on and turn-off dv/dt and di/dt figures use the maximum absolute slope
over a 1 ns interval. Raw sample derivatives are also available in the CSV.
Reported maxima depend on time resolution and bandwidth. Earlier nominal
timestep/window convergence checks are retained in results/double\_pulse;
they are not a physical validation of every corner.

\subsection{How to read the figures}
Rows compare bus voltage; columns compare bus-side loop inductance; colors
compare external gate resistance. Lines connect sampled points for readability,
without implying a validated interpolation model. Red crosses mark cases failing
the 90 V design ceiling or package gate limits (-4 V, +6 V); both devices are
checked. The voltage figure also marks the 100 V drain absolute rating.
Failed points remain visible for sensitivity analysis and must not be treated
as qualified loss data. The model continues above ratings because it has no
destructive breakdown model. Actual PCB parasitics, gate driver, dynamic
on-resistance and production spread still require validation.
@FIGURES@
\end{document}
